\documentclass[11pt,a4paper]{article}
\usepackage[margin=27mm]{geometry}
\usepackage{amsmath,amssymb,amsthm,fontspec,unicode-math,booktabs,longtable,hyperref,fvextra}
\setmainfont{STIXGeneral}\setmathfont{STIX Two Math}\setmonofont{Menlo}[Scale=MatchLowercase]
\setlength{\emergencystretch}{3em}
\newtheorem{theorem}{Theorem}\newtheorem{lemma}[theorem]{Lemma}\newtheorem{corollary}[theorem]{Corollary}\newtheorem{proposition}[theorem]{Proposition}
\newcommand{\Aut}{\operatorname{Aut}}\newcommand{\Orb}{\operatorname{Orb}}\newcommand{\supp}{\operatorname{supp}}
\hypersetup{colorlinks=true,linkcolor=blue,urlcolor=blue}

\title{Joins, wreath products and reversal of polytope markings}
\author{Human Compression Project}\date{}
\begin{document}\maketitle
\begin{abstract}
We apply the common orbit construction to joins, permutation wreath products
and reversal of regular polytope markings. Coordinate formulas give the full
parabolics and the exact panel sizes, including nontrivial chamber
stabilizers. An explicit module automorphism identifies the three pairs of
signed permutation constructions in the source. Infinite joins lead to a
direct description of the restricted and unrestricted lamplighter geometries.
The application records corrections to several unrestricted source statements.
The ten affirmative sporadic cases have separate finite certificates [5,6].
Their finite input data and computational proof obligations are recorded
there rather than being assumed in the orbit argument.
\end{abstract}

\section{Scope and common input}
The original is Piedade and Tranchida's \emph{The geometry of wreath and
semi-direct products}, arXiv:2605.04646v2, twenty-two pages [1]. We use
the complete companion paper \emph{Orbit intersections and marked semidirect
products} [2]. Its orbit construction permits nontrivial Borel subgroups;
its faithful regular specialization identifies the full group of
type-preserving automorphisms. Its module and cocycle theorems will be used
by name. None of those proofs is assumed from the original.

For a coset geometry $\alpha=(A,(A_j)_{j\in S})$ write
$A_E=\bigcap_{j\in E}A_j$, with $A_\varnothing=A$. The Borel subgroup
is $A_S$. For a faithful regular hypertope, distinguish the involutions
$a_j$ associated with the base chamber and put
$A(R)=\langle a_j:j\in R\rangle$. Then $A_E=A(S\setminus E)$.
Keeping these two notations separate is essential when a Borel subgroup
is nontrivial. All cosets below are left cosets.

The source's preliminary Lemma 2.4 identifies certain generated
parabolics with complementary intersections without requiring a trivial
Borel subgroup. Its empty generated subgroup is $1$, so these equalities
cannot hold at the extremes unless the Borel is trivial. A separate
generation hypothesis is needed for the converse intersection criterion.
For example, in $G=\langle x,y,z\rangle=C_2^3$, the parabolics
$\langle y,z\rangle,\langle x,z\rangle$ give the connected geometry
$K_{2,2}$ but have intersection $\langle z\rangle\ne1$. Using instead
$\langle x\rangle,\langle y\rangle$, the four generated-support
subgroups $1,\langle x\rangle,\langle y\rangle,\langle x,y\rangle$
satisfy every support-intersection identity, although the coset geometry
is two disjoint copies of $K_{2,2}$. The complete common paper uses
residual generation with the Borel retained, which applies to both faithful
and unfaithful models.

\section{Joins and their two group models}
Let $(\Gamma_\omega)_{\omega\in\Omega}$ be coset incidence systems,
with groups $H^\omega$ and finite type sets $I_\omega$. Their type sets
are regarded as disjoint. The join has their disjoint union of elements,
retains incidence within each factor, and declares every pair in distinct
factors incident. We allow $\Omega$ to be infinite in this section.

\begin{proposition}\label{prop:joins}
Both the direct product $\prod_\omega H^\omega$ and the restricted
direct product $\bigoplus_\omega H^\omega$ realize the join as a coset
incidence system. For a type $i\in I_\lambda$, its maximal parabolic
restricts the $\lambda$ coordinate to $H_i^\lambda$ and leaves every
other coordinate unrestricted, within the chosen product.

If every factor action is flag-transitive, the full product action is
flag-transitive. The restricted product carries one flag to another of
the same type whenever they differ in only finitely many factors. Every
join of geometries is a geometry; it is residually connected, firm or
thin when all factors have the corresponding property.
\end{proposition}
\begin{proof}
Projection to coordinate $\lambda$ induces a bijection from the cosets
of the displayed maximal parabolic to $H^\lambda/H_i^\lambda$.
Surjectivity holds for the restricted product as well: a representative
can have just its $\lambda$ coordinate nonidentity. Two cosets with
types in the same factor intersect exactly when the projected cosets
intersect. For two different factors the coordinate requirements can
always be satisfied simultaneously by an element with at most two
nonidentity coordinates. These statements identify both incidence
systems with the join, without any finiteness assumption on $\Omega$.

A flag is the union of its component flags. Extend each component to
a chamber; their union is a chamber of the join. Conversely a maximal
flag must be maximal in each component. Choose, in each coordinate,
an element carrying one component flag to the other. Their tuple lies
in the full product. If only finitely many component flags differ, take
the identity in every other coordinate to obtain a restricted-product
element. These choices use the usual axiom of choice for arbitrary
families.

A residue is the join of the corresponding factor residues, after
discarding those with no missing types. If at least two nonempty factors
remain, the incidence graph is connected because any two vertices in
one factor are connected through a vertex of another factor. If exactly
one factor remains and the residue has rank at least two, connectivity
is the residual connectivity of that factor. Finally a panel misses
one type in one factor, so its number of completions is exactly that of
the corresponding factor panel. This proves the last assertions.
\end{proof}

For two finite-rank factors the group is $A\times B$, and direct
intersection gives
\[
 G_J=A_{J\cap S}\times B_{J\cap I}.
\]
Its geometry is finite exactly when both factors are finite. These
observations recover source Lemma 3.2 and Theorems 3.3, 3.6 and 3.7.
The source's assertion that the restricted product is \emph{never}
flag-transitive for infinite $\Omega$ needs a qualification. It fails
to be flag-transitive if infinitely many factors permit distinct flags
of the same type: choose such flags in every one of those factors, and
no finitely supported element can carry their unions to each other.
An infinite family of one-chamber factors is a counterexample to the
unqualified assertion. In the firm examples used below every nonempty
factor does permit different chambers, so the stated obstruction applies.

\section{Wreath products and exact panel sizes}
Let $\alpha=(A,(A_j)_{j\in S})$ and
$\beta=(B,(B_i)_{i\in I})$ be flag-transitive geometries of finite,
nonzero rank. Let $\Omega$ be a finite nonempty $B$-set. For every orbit
$O\subseteq\Omega$ and every $j\in S$, choose a point
$\lambda_{O,j}\in O$ satisfying
\begin{equation}\label{eq:wreath-ipo}
 B_H\lambda_{O,j}\cap B_K\lambda_{O,j}
       =B_{H\cup K}\lambda_{O,j}\qquad(H,K\subseteq I).
\end{equation}
The choices may depend on $j$. Let $D=A^\Omega$, with
$b(v)_\omega=v_{b^{-1}\omega}$, so $D\rtimes B=A\wr_\Omega B$.
The type set of the join $\alpha^\Omega$ is $S\times\Omega$.
Its orbit types are $L_{O,j}=\{j\}\times O$.

\begin{proposition}\label{prop:wreath}
The coordinate action is admissible for the common orbit construction
exactly when the selected representatives satisfy
\eqref{eq:wreath-ipo}. The wreath geometry has the following maximal
parabolics. For $i\in I$ put
\[
 E_i(\omega)=\{j\in S:\omega\notin B_i\lambda_{O(\omega),j}\}.
\]
Then
\begin{align}
 G_i&=\left(\prod_{\omega\in\Omega}A_{E_i(\omega)}\right)\rtimes B_i,
          &&i\in I,\label{eq:wreath-old}\\
 G_{L_{O,j}}&=\left(\prod_{\omega\in O}A_j\times
                     \prod_{\omega\notin O}A\right)\rtimes B.
          \label{eq:wreath-new}
\end{align}
It is a flag-transitive geometry, residually connected and firm if the
inputs are so. Its old panel sizes are
$[B_{I\setminus\{i\}}:B_I]$, and its new panel size at $L_{O,j}$ is
\begin{equation}\label{eq:wreath-panel}
 [A_{S\setminus\{j\}}:A_S]^{\,|B_I\lambda_{O,j}|}.
\end{equation}
In particular, for thin inputs it is thin exactly when $B_I$ fixes
all the selected representatives.
\end{proposition}
\begin{proof}
The action carries a coordinate parabolic indexed by $(j,\omega)$
to that indexed by $(j,b\omega)$. Its orbit is $L_{O,j}$, with the
same orbit intersection identity as the action on $O$. This proves
admissibility and its converse with the indexed action understood.

The common theorem takes the intersection of coordinate parabolics
outside the union of the orbits $B_i(j,\lambda_{O,j})$. At coordinate
$\omega$ the retained restrictions are exactly the $j\in E_i(\omega)$,
which proves \eqref{eq:wreath-old}. Intersecting every coordinate
parabolic in $L_{O,j}$ gives \eqref{eq:wreath-new}. The common theorem
and Proposition~\ref{prop:joins} give flag transitivity, residual
connectivity and firmness.

The Borel subgroup of the product geometry is $A_S^\Omega$. Deleting
the new type $L_{O,j}$ relaxes the restriction from $A_S$ to
$A_{S\setminus\{j\}}$ precisely at the coordinates
$B_I\lambda_{O,j}$. Taking the index over $A_S^\Omega$ gives
\eqref{eq:wreath-panel}. For thin inputs its base is two, so the
index is two if and only if that nonempty orbit has one point.
The old panels already have size two. This proves the equivalence.
\end{proof}

If representatives are chosen coherently, so $\lambda_{O,j}=\lambda_O$
for every $j$, the coordinate factor in \eqref{eq:wreath-old} is $A$
on $\bigcup_O B_i\lambda_O$ and $A_S$ outside it. The printed source
Proposition 4.3 reverses this support and omits the Borel factors. The
formula above follows directly from its defining intersections and also
covers independent choices for different types.

Source Theorem 4.4 asserts thinness without the necessary orbit condition.
Take $\alpha=(C_2,1)$, $\beta=(S_3,C_3)$ of rank one and let $S_3$
act on three coordinates. Both geometries have two elements and are
thin. The orbit identity is automatic for rank one. But
$|C_3\lambda|=3$, so the new panel has size $2^3=8$. More explicitly,
the two maximal parabolics are $C_2^3\rtimes C_3$ and $S_3$ in a
group of order $48$. Their coset geometry is the complete bipartite
graph $K_{2,8}$: the product of the parabolics is the whole group.
It is firm and connected, but not thin. Passing from $S_3$ to $S_3/C_3$
does not repair this construction: the original action on three coordinates
does not factor through that quotient.

For faithful regular inputs the Borel groups are trivial. Every permitted
wreath construction is then a regular hypertope, with full type-preserving
automorphism group $A\wr_\Omega B$. This is the precise regular case
used in the source applications.

\section{The six signed permutation families}
Let $r\geq2$ and represent $\mathbb F_2^r\rtimes S_r$ on $2r$ points
in $r$ blocks of size two. A vector $e_j$ flips block $j$ and the
adjacent transposition $s_i=(i\ i+1)$ interchanges two adjacent blocks
without internal flips. Put $\lambda=\sum e_j$. The complete common
module theorem applies to each seed $\lambda,e_1,\lambda+e_1$;
the first mark is its block-flip vector and the remaining marks are
$s_1,\ldots,s_{r-1}$.

\begin{corollary}\label{cor:six}
All six lists obtained by using one of these three seeds and either
the marks $(0,s_i)$ or the marks $(\lambda+e_i+e_{i+1},s_i)$ are
string C-groups. They give the diagrams (1), (2) of Table 9 and
(13)--(16) of Table 6 in the bound Fernandes--Piedade catalogue [3].
Their abstract groups, in the order of the three seeds, are
$C_2\times S_r$, $C_2\wr S_r$, and $C_2\wr S_r$ for even $r$ or
$\mathbb F_2^{r,\mathrm{even}}\rtimes S_r$ for odd $r$.
\end{corollary}
\begin{proof}
The ordinary markings satisfy every subgroup intersection by the common
module theorem. Its change of coordinates fixes the first mark and sends
each $(0,s_i)$ to $(\lambda+e_i+e_{i+1},s_i)$. That latter vector
means exactly the vertical $i$-edges in every block except $i$ and
$i+1$ in the source diagrams. Thus this is an identification of the
actual marked permutations, not merely of their abstract groups. The
three vector spans and their sole odd-rank relation are calculated in
the common proposition on signed markings. Transport through its
automorphism proves the assertion.
\end{proof}

The auxiliary classification in [3] has range $r=n/2\geq7$; the
algebraic construction just proved works for every $r\geq2$. Source
Proposition 5.1 covers five of these markings in both parities and the
sixth only for even $r$. The preceding paragraph of that source leaves
odd-rank graph (16) unverified. Corollary~\ref{cor:six} proves its
intersection property too. This is a solution of that explicitly recorded
gap; it is not a historical-priority assertion. It does not prove the
external classification of all possible high-rank permutation graphs.

\section{Reversal of a regular polytope}
Let $\alpha$ be a self-dual regular polytope of rank $n\geq3$, with
distinguished involutions $a_0,\ldots,a_{n-1}$ and group $A$. A duality
reverses types. Composing a duality with an automorphism of $\alpha$
makes it preserve the base chamber as a set. Its square is type-preserving
and fixes that chamber, so is the identity by regularity. Conjugation
therefore gives an involution $\tau\in\Aut(A)$ satisfying
\begin{equation}\label{eq:reversal}
 \tau(a_i)=a_{n-1-i}.
\end{equation}
It is nonidentity because the distinguished involutions are distinct.

\begin{proposition}\label{prop:fold}
Choose one representative from each orbit of the reversal
$i\mapsto n-1-i$. Together with the outer factor $t$ of order two,
the selected $a_i$ are the distinguished involutions of a regular
hypertope for $A\rtimes\langle t\rangle$, where $t$ acts by $\tau$.
Its rank is $\lceil n/2\rceil+1$. For selected $i,j$ their mutual
label is $\operatorname{ord}(a_ia_j)$. Their labels with $t$ are
\[
 \operatorname{ord}(a_it)=
 \begin{cases}
 2,&i=n-1-i,\\
 2\operatorname{ord}(a_i a_{n-1-i}),&i\ne n-1-i.
 \end{cases}
\]
For $n=3,4,5$ representatives can always be chosen so that the marking
is a string C-group. For $n\geq6$ its diagram is non-linear for every
choice of representatives. For $n=4,5$, a non-linear choice exists when
$\operatorname{ord}(a_0a_1)>2$.
\end{proposition}
\begin{proof}
Use the rank-one geometry $(\langle t\rangle,1)$ in the faithful
regular orbit construction. For a rank-one group, the only parabolics
are the whole group and the identity, and their orbit intersection
identity holds at every point. There are $\lceil n/2\rceil$ reversal
orbits. The common theorem gives the asserted geometry, generation and
full automorphism group. Its exact doubled-label lemma gives the displayed
order for a nonfixed index. At a fixed index $t$ commutes with $a_i$;
their product is a nonidentity involution since it projects to $t$.

For $n=3$ choose $0,1$. The mark $a_1$ commutes with $t$, giving
the string order $a_1,a_0,t$. For $n=4$ choose $0,2$. The old marks
commute, giving the order $a_0,t,a_2$. The source's printed choice
$0,3$ selects the same reversal orbit twice; $0,2$ is the required
correction. For $n=5$ choose $0,2,3$, with string order
$a_0,t,a_3,a_2$. Every nonconsecutive pair commutes, by the displayed
orders and the original string relations. The ordinary correspondence
between string C-groups and regular polytopes gives the polytope claims,
including possible labels two.

Every nonfixed reversal orbit contributes a neighbour of $t$ in the
diagram. Indeed $a_i$ and $a_{n-1-i}$ are distinct involutions, so their
product has order at least two and its doubled label is at least four.
There are $\lfloor n/2\rfloor$ such orbits. For $n\geq6$ the degree
of $t$ is therefore at least three. For $n=4,5$ select $0,1$ and also
the fixed middle index when present. If $\operatorname{ord}(a_0a_1)>2$,
the three nodes $a_0,a_1,t$ form a triangle. In rank five the full
diagram is a triangle with an additional edge to the middle mark when
that label exceeds two; it is not the square described in the source.
\end{proof}

The unconditional non-linear assertion of source Theorem 5.2(b) is false
in rank four. To specify an admissible all-digon polytope of rank $n$,
take two faces at each proper rank and one least and greatest face.
Every face of lower rank is below every face of higher rank. Intervals
of rank difference two have two intermediate faces. Chambers are
independent binary choices at the $n$ proper ranks. Changing one choice
at a time proves strong flag connectivity, including while retaining
any fixed common faces. Thus this poset is an abstract regular polytope.
Its full automorphism group independently swaps the two faces at each
rank, so is $C_2^n$, with the standard basis as marks. Rank reversal
is a duality fixing the chamber of zero choices.

Now take $n=4$. In $A\rtimes\langle t\rangle$, all
involutions in $A$ commute. An involution outside $A$ has form $v t$
with $\tau(v)=v$; any two such involutions commute as well. Every
triple of involutions therefore contains a commuting pair. In particular
no rank-three hypertope with triangular diagram has this group. The
common construction and the polytope conclusion remain valid; the lost
claim is the unrestricted non-linear choice. Indeed every regular
hypertope has distinguished involutions, and a commuting pair gives a
rank-two residue of gonality two, hence no edge between those types.
A non-linear diagram on three nodes would require all three edges.

\begin{lemma}\label{lem:automorphism-map}
Suppose $\tau$ is not inner. The natural map
\[
 A\rtimes\langle t\rangle\longrightarrow\Aut(A),\qquad
 (a,t^\epsilon)\longmapsto\operatorname{Inn}(a)\tau^\epsilon
\]
has kernel $Z(A)\times1$. It is injective exactly when $A$ is centerless.
\end{lemma}
\begin{proof}
Conjugation intertwines the action of $\tau$, so the displayed map is
a homomorphism. A kernel element with $\epsilon=1$ would express
$\tau$ as an inner automorphism. With $\epsilon=0$ its image is the
identity exactly when $a$ is central. This proves both statements.
\end{proof}
Source Theorem 5.2(d) omits the center. For the rank-three all-digon
polytope the reversal is outer, the semidirect group has order $16$, and
$|\Aut(C_2^3)|=|\mathrm{GL}_3(2)|=7\cdot6\cdot4=168$.
Lagrange's theorem rules out even an abstract embedding. For sporadic
simple groups the centerless hypothesis does hold.

Nor does $\operatorname{Out}(A)$ automatically act on $A$: it is a
quotient of $\Aut(A)$, and a lift as a group action is additional data.
For $A=D_{16}$ the quotient map $\Aut(A)\to\operatorname{Out}(A)$
has no section. Here is a direct proof, also obtained independently in [4].
Write $A=\langle r,s\mid r^8=s^2=1,srs=r^{-1}\rangle$. Every
automorphism has form $f_{a,b}:r\mapsto r^a,s\mapsto r^bs$,
where $a\in\{1,3,5,7\}$ and $b\in\mathbb Z/8$. This follows
because $\langle r\rangle$ is the unique cyclic subgroup of order
eight; all these substitutions do give automorphisms. Composition is
$f_{a,b}f_{c,d}=f_{ac,b+ad}$. The inner subgroup has $a=1,7$ and
even $b$, so the outer quotient is $C_2^2$, generated by the classes
$X=[f_{1,1}]$ and $Y=[f_{3,0}]$.
An involutory lift of $X$ is $f_{7,c}$ with $c$ odd. An involutory
lift of $Y$ is either $f_{3,b}$ with $b$ even, or $f_{5,b}$ with
$b=0,4$. Commutation with the first kind requires $c-3b=0$ modulo
four, impossible by parity; with the second kind it requires
$4c-6b=0$ modulo eight, also impossible. A section would provide
commuting involutory lifts of $X,Y$, a contradiction. Throughout this
application we specify actual automorphisms such as \eqref{eq:reversal}.

\section{Two markings of the seven-generator Coxeter example}
In source Example 3.9 the Coxeter diagram is a three-armed tree with
edges $01,12,03,34,05,56$, all labelled three. All other pairs commute.
Let $b_7$ interchange arms $(1,2)$ and $(3,4)$ and let $b_8$
interchange arms $(3,4)$ and $(5,6)$. They generate $S_3$, acting by
Coxeter diagram automorphisms. Its type orbits are
$\{0\}$, $\{1,3,5\}$ and $\{2,4,6\}$.

For the natural action of $S_3$ on three points, its two rank-one
parabolics are the subgroups generated by the two adjacent transpositions.
At any selected point, their orbits intersect only in that point: at an
endpoint one orbit is a singleton, and at the middle point the two
two-element orbits have only the middle point in common. This proves
the full orbit identity, since the other parabolics are $1$ and $S_3$.
Thus any choice of one representative from each orbit is admissible.

For representatives $0,1,2$, the old maximal parabolics are
\[
 G_7=A_{\{3,4,5,6\}}\rtimes\langle b_8\rangle,
 \qquad G_8=A_{\{5,6\}}\rtimes\langle b_7\rangle,
\]
and the three orbit parabolics are $A_L\rtimes S_3$. The resulting
five marks have the following noncommuting pairs:
\[
 a_0a_1:3,\quad a_1a_2:3,\quad
 a_1b_7:4,\quad a_2b_7:4,\quad b_7b_8:3.
\]
Every omitted pair commutes. The doubled labels follow because $a_1$
commutes with $a_3$ and $a_2$ with $a_4$ in the Coxeter group.
For representatives $0,1,6$, the five marks instead form the path
\[
 a_0\ \mathrel{\overset{3}{-}}\ a_1\
 \mathrel{\overset{4}{-}}\ b_7\
 \mathrel{\overset{3}{-}}\ b_8\
 \mathrel{\overset{4}{-}}\ a_6.
\]
The common regular theorem proves the complete intersection and geometry
properties in both cases, rather than just the pair relations. In the
second case the string correspondence supplies a regular polytope.

The parabolic change noted in the source is also exact. Replacing
representative $2$ by $4$ changes $G_7$ to
$A_{\{2,3,5\}}\rtimes\langle b_8\rangle$. The first $G_7$ is
$S_4\times C_2$, since the support $\{0,1,2\}$ is an $A_3$ path
fixed by $b_8$. The changed $G_7$ is $S_3\times D_8$: its support
$\{0,1,4,6\}$ has an $A_2$ edge and two isolated vertices, and
$b_8$ swaps the latter. They are nonisomorphic, as their abelianizations
have orders four and eight. The Coxeter parabolic embedding theorem is
the retained classical input in these identifications.

\section{The tetrahedron and the cube}
For the tetrahedron take $A=S_4$, $a_0=(1\ 2)$, $a_1=(2\ 3)$,
$a_2=(3\ 4)$, and let $t$ act by conjugation by $(1\ 4)(2\ 3)$.
The representatives $0,1$ give marks $a_1,a_0,t$ of string type
$\{3,4\}$ (the dual ordering has type $\{4,3\}$). The resulting
group has order $48$. Its maximal parabolics have orders $4,8,6$,
and therefore its three element counts are $12,6,8$ in the original
orbit/old-type order.

Here is an incidence identification independent of a group-order name.
The eight cube vertices are the four tetrahedron vertices and its four
faces. An edge is an incident pair $(v,F)$ with $v$ a vertex of the
face $F$; a cube face is a tetrahedron edge $e$. Its boundary comprises
the two endpoints of $e$ and the two faces containing $e$, with all
four corresponding incident pairs. This is a quadrilateral. Put the
tetrahedron at the four cube vertices whose sign-coordinate product is
$+1$. Associate a tetrahedron face to the antipode of its omitted
vertex. A tetrahedron vertex lies in that face exactly when the two
associated cube vertices differ in one coordinate. The six tetrahedron
edges yield the six coordinate faces of the cube by the preceding rule.
These bijections preserve every incidence, proving the cube identification
in source Example 3.18. Exchanging point and face types gives its dual
polytope ordering.

\section{The two lamplighter geometries}
Let $D$ be either $D_{\mathrm{fin}}=\bigoplus_{\mathbb Z}C_2$ or
$D_{\mathrm{all}}=\prod_{\mathbb Z}C_2$, written as binary lamp
configurations. Let $t$ shift coordinates and let $e_n$ toggle lamp $n$.
In $G=D\rtimes\langle t\rangle$ take maximal parabolics
$P=\langle t\rangle$ and $Q=\langle e_0\rangle$.

\begin{proposition}\label{prop:lamps}
The coset geometry $(G;(P,Q))$ is the subdivision of the graph whose
vertices are configurations in $D$ and whose edges join configurations
differing at exactly one lamp. It is flag-transitive. Its connected
components are indexed by $D/D_{\mathrm{fin}}$. In particular it is
connected for restricted configurations and disconnected for unrestricted
configurations. Every edge-type element is incident with two configuration
elements; every configuration is incident with countably infinitely many
edge-type elements.
\end{proposition}
\begin{proof}
Write a group element as $(v,n)$, with product
$(v,n)(w,m)=(v+t^nw,n+m)$. Its $P$-coset is determined by $v$,
giving the configuration vertices. Its $Q$-coset consists of
$(v,n)$ and $(v+e_n,n)$; it records a configuration with its $n$th
lamp left unspecified. Incidence with a $P$-coset means resolving
that lamp in one of its two ways. This proves the description and all
valencies directly.

Any incident pair of cosets has a common representative. Left
translation takes the standard incident pair to that pair, proving
flag transitivity in rank two. A path changes finitely many lamps.
Conversely, if two configurations differ on a finite set, toggle those
lamps in turn to obtain a path. Components are consequently precisely
the cosets of $D_{\mathrm{fin}}$ in $D$. Each unspecified-lamp element
lies in the component of its two resolutions. Finally unrestricted
configurations can differ at infinitely many lamps, so there is more
than one component in that case.
\end{proof}

For completeness the restricted group has presentation
\[
 \langle a,t\mid a^2=1,\ (a t^n a t^{-n})^2=1\ (n\geq1)\rangle.
\]
Indeed the conjugates $a_n=t^nat^{-n}$ are involutions. The displayed
relations and their conjugates make every two of them commute, so
every word has the form $\prod_{n\in F}a_n\,t^k$ with $F$ finite.
The concrete semidirect product satisfies the presentation, and in it
that normal form is the identity only for $F=\varnothing,k=0$.
The mutually inverse maps therefore give the claimed presentation.
The source's displayed presentation on page 20 omits the square on
$a t^n a t^{-n}$. Read literally, its $n=1$ relator makes $a$ commute
with $t$, and the presented group is $C_2\times\mathbb Z$.
The squared relators above correct that auxiliary presentation.
We do not need or reprove the separate classical theorem that the
restricted lamplighter group is not finitely presentable.

This description recovers source Section 5.3 and pinpoints the failure
of the finite-type residual-connectivity theorem for unrestricted
infinite joins. The group generated by $P$ and $Q$ is only
$D_{\mathrm{fin}}\rtimes\mathbb Z$, a proper subgroup of
$D_{\mathrm{all}}\rtimes\mathbb Z$. The missing hypothesis cannot
be removed by passing to the full direct product.

\section{Attribution and certificates}
The constructions, examples and requested results are those of [1];
the six permutation diagrams come from [3]. The join description and
the definition of generalized twisting are retained constructions.
The common proof [2] supplies the complete orbit argument, the module
intersection proof and its explicit change of coordinates. The present
paper supplies their coordinate and geometric applications. The
correction note [4] gives elementary countermodels for the unqualified
thinness, non-linear and automorphism-embedding claims. These are
corrections to assertions, not newly solved open conjectures.

The generic cocycle automorphism and transport of subgroup intersections
have scoped Lean proofs. They do not formalize the complete geometric
arguments, the source diagram identifications or the classification of
sporadic groups. Finite validation of signed markings is supplementary;
the arbitrary-rank assertion is the written module theorem. The ten
affirmative sporadic cases in source Theorem 5.3 are supplied by [5,6],
including the Harada--Norton case absent from the source proof. Each
certificate establishes both a rank-three regular polytope and a
non-linear rank-three regular hypertope. These certificates retain their
named group data and, where used, the specified maximal-subgroup
classification; those inputs are not consequences of the orbit theorem.
The source exceptions concerning $\mathrm{Fi}_{22}$ and
$\mathrm{Fi}_{24}^{\prime}$ are outside this paper's recovery contract.
No length factor or claim of complete formalization follows from the
number of cases.

\begin{thebibliography}{9}
\bibitem{source} C. A. Piedade and P. Tranchida, \emph{The geometry of wreath and semi-direct products}, arXiv:2605.04646v2, 2026.
\bibitem{common} Human Compression Project, \emph{Orbit intersections and marked semidirect products}, companion proof edition.
\bibitem{graphs} M. E. Fernandes and C. A. Piedade, \emph{Regular polytopes of rank $n/2$ for transitive groups of degree $n$}, arXiv:2407.16003v2, 2025; J. Algebra 689 (2026), 256--296.
\bibitem{corrections} Human Compression Project, \emph{Borel subgroups and duality in coset twisting}, independently prepared five-page proof supplement. Its correction proofs used here are included in full above.
\bibitem{sporadic} Human Compression Project, \emph{Finite witnesses for nine sporadic extensions}, companion certificate packet.
\bibitem{hn} Human Compression Project, \emph{A finite witness for the Harada--Norton extension}, companion certificate packet.
\end{thebibliography}
\end{document}
